
\documentclass{article}
\usepackage{colortbl}
\usepackage{makecell}
\usepackage{multirow}
\usepackage{supertabular}

\begin{document}

\newcounter{utterance}

\centering \large Interaction Transcript for game `eqbench', experiment `eq\_bench', episode 138 with reflexion{-}Qwen3.5{-}27B.
\vspace{24pt}

{ \footnotesize  \setcounter{utterance}{1}
\setlength{\tabcolsep}{0pt}
\begin{supertabular}{c@{$\;$}|p{.15\linewidth}@{}p{.15\linewidth}p{.15\linewidth}p{.15\linewidth}p{.15\linewidth}p{.15\linewidth}}
    \# & \multicolumn{2}{c}{Player} && \multicolumn{2}{c}{Game Master} \\
    \hline

    \theutterance \stepcounter{utterance}  
    & & & \multicolumn{4}{p{0.6\linewidth}}{
        \cellcolor[rgb]{0.9,0.9,0.9}{
            \makecell[{{p{\linewidth}}}]{
                \texttt{\tiny{[P1$\langle$GM]}}
                \texttt{You have attempted this same task before and did not succeed. These are the plans you wrote after those attempts.} \\
\texttt{Trial \#1: New plan: My previous attempt failed because I assigned a high score to "Offended" (7) and "Indignant" (6) without sufficiently distinguishing the nuance of Henrietta's response. Her reply, "My home does not need any more statements... It speaks volumes of my life as it is," is a firm, polite rejection of unsolicited advice, not an expression of anger or insult. I incorrectly interpreted her boundary{-}setting as hostility. Next time, I will analyze the tone of the rejection more carefully: since she remains polite ("Mr. Harrington") and focuses on her own values rather than attacking his character, "Offended" and "Indignant" should be scored very low (0{-}2). I will instead look for subtle cues of "Anxious" (perhaps regarding the intrusion) or simply accept that the primary emotion is a neutral "Satisfied" with her own life, resulting in a score of 0 for the negative emotions and a low score for any potential mild annoyance, ensuring the scores reflect a polite but firm dismissal rather than an emotional outburst.} \\
\\ 
\texttt{Trial \#2: New plan: My previous attempt failed because I assigned a high score (8) to "Satisfied," misinterpreting Henrietta's polite but firm rejection as contentment. Her statement, "My home does not need any more statements," is a defensive boundary against an unsolicited sales pitch, not an expression of happiness. I incorrectly assumed her politeness equated to satisfaction. Next time, I will analyze the subtext of her refusal: she is rejecting an intrusion on her personal space and values. Therefore, "Satisfied" should be scored 0. Instead, I will assign higher scores to "Indignant" (reflecting her annoyance at the presumption that she needs his "statement piece") and "Offended" (reflecting the slight insult to her current home's sufficiency), while keeping "Anxious" low as there is no threat of danger, only social friction. I will ensure the highest score reflects the negative reaction to the unsolicited advice.} \\
\\ 
\texttt{Trial \#3: New plan: My previous attempts failed because I misinterpreted the nuance of Henrietta's polite but firm rejection. In the last attempt, I scored "Indignant" too high (6) and "Offended" too low (3), failing to recognize that her response ("My home does not need any more statements... It speaks volumes of my life as it is") is a dignified assertion of self{-}worth rather than an angry outburst. She is not "indignant" (which implies righteous anger) nor deeply "offended" (which implies personal insult), but rather mildly annoyed or dismissive of his presumption. The key mistake was overestimating the intensity of the negative emotions. Next time, I will assign a low score (1{-}2) to "Indignant" to reflect mild annoyance at the unsolicited advice, a very low score (0{-}1) to "Offended" as she remains polite, and 0 to "Satisfied" and "Anxious" as they do not fit the context. The strategy will be to treat her response as a calm boundary{-}setting, resulting in low scores across the board, with "Indignant" being the only non{-}zero score to capture the slight friction of the rejection.} \\
\\ 
\\ 
\\ 
\texttt{Your task is to predict the likely emotional responses of a character in this dialogue:} \\
\\ 
\texttt{Barty: Ah, Mrs. Adams, how delightful to see you. I trust you are well?} \\
\texttt{Henrietta: Mr. Harrington, I'm as well as can be expected. What brings you here today?} \\
\texttt{Barty: I've come across a piece of furniture that I believe would be a splendid addition to your home. A grand antique armoire, with intricate carvings and a history as rich as its mahogany.} \\
\texttt{Henrietta: And why, pray tell, would I need an armoire, Mr. Harrington?} \\
\texttt{Barty: Well, Mrs. Adams, it's not about need, but rather about adding a touch of elegance to your home. It's a statement piece, a testament to your refined taste.} \\
\texttt{Henrietta: My home does not need any more statements, Mr. Harrington. It speaks volumes of my life as it is.} \\
\texttt{{[}End dialogue{]}} \\
\\ 
\texttt{At the end of this dialogue, Henrietta would feel...} \\
\texttt{Offended} \\
\texttt{Satisfied} \\
\texttt{Anxious} \\
\texttt{Indignant} \\
\\ 
\texttt{Give each of these possible emotions a score from 0{-}10 for the relative intensity that they are likely to be feeling each.} \\
\\ 
\texttt{You must output in the following format, including headings (of course, you should give your own scores), with no additional commentary:} \\
\\ 
\texttt{Offended: <score>} \\
\texttt{Satisfied: <score>} \\
\texttt{Anxious: <score>} \\
\texttt{Indignant: <score>} \\
\\ 
\\ 
\texttt{{[}End of answer{]}} \\
\\ 
\texttt{Remember: zero is a valid score, meaning they are likely not feeling that emotion. You must score at least one emotion > 0.} \\
\\ 
\texttt{Your answer:} \\
            }
        }
    }
     \\ \\

    \theutterance \stepcounter{utterance}  
    & \multicolumn{4}{p{0.6\linewidth}}{
        \cellcolor[rgb]{0.561,0.737,0.561}{
            \makecell[{{p{\linewidth}}}]{
                \texttt{\tiny{[P1$\rangle$GM]}}
                \texttt{Offended: 2} \\
\texttt{Satisfied: 0} \\
\texttt{Anxious: 0} \\
\texttt{Indignant: 3} \\
            }
        }
    }
    & & \\ \\

    \theutterance \stepcounter{utterance}  
    & & & \multicolumn{2}{p{0.3\linewidth}}{
        \cellcolor[rgb]{0.9,0.9,0.9}{
            \makecell[{{p{\linewidth}}}]{
                \texttt{\tiny{[GM$|$GM]}}
                \texttt{Offended: 2} \\
\texttt{Satisfied: 0} \\
\texttt{Anxious: 0} \\
\texttt{Indignant: 3} \\
            }
        }
    }
    & & \\ \\

    \theutterance \stepcounter{utterance}  
    & & & \multicolumn{2}{p{0.3\linewidth}}{
        \cellcolor[rgb]{0.9,0.9,0.9}{
            \makecell[{{p{\linewidth}}}]{
                \texttt{\tiny{[GM$|$GM]}}
                \texttt{Offended: 6} \\
\texttt{Satisfied: 2} \\
\texttt{Anxious: 0} \\
\texttt{Indignant: 8} \\
            }
        }
    }
    & & \\ \\

    \theutterance \stepcounter{utterance}  
    & & & \multicolumn{2}{p{0.3\linewidth}}{
        \cellcolor[rgb]{0.9,0.9,0.9}{
            \makecell[{{p{\linewidth}}}]{
                \texttt{\tiny{[GM$|$GM]}}
                \texttt{game\_result = LOSE} \\
            }
        }
    }
    & & \\ \\

    \theutterance \stepcounter{utterance}  
    & & \multicolumn{4}{p{0.6\linewidth}}{
        \cellcolor[rgb]{0.95,0.95,0.95}{
            \makecell[{{p{\linewidth}}}]{
                \texttt{\tiny{[P1$\leftrightarrow$P]}}
                \texttt{Post{-}episode Reflexion (not sent to GM):} \\
\texttt{New plan: My previous attempts failed because I oscillated between overestimating the intensity of negative emotions (Indignant/Offended) and misinterpreting the context as positive (Satisfied), failing to pinpoint the specific nuance of a polite but firm social rejection. In the last attempt, I scored "Indignant" at 3 and "Offended" at 2, which might still be too high for a character who maintains formal address ("Mr. Harrington") and uses a metaphorical, non{-}aggressive defense ("speaks volumes"). The core mistake was assigning significant intensity to emotions that imply active anger or deep personal insult, whereas Henrietta's response is a calm, dignified dismissal. Next time, I will interpret her tone as "mildly annoyed" rather than "indignant" or "offended." I will assign a score of 1 to "Indignant" to capture the slight friction of the unsolicited pitch, and 0 to "Offended," "Satisfied," and "Anxious," ensuring the scores reflect a low{-}intensity, polite boundary setting rather than an emotional outburst or contentment.} \\
            }
        }
    }
    & \\ \\

\end{supertabular}
}

\end{document}
